Due to the complex and unpredictable nature of vehicle control and environmental perception, autonomous vehicle navigation in unstructured environments is challenging~\cite{borges2022survey}. Learning-based navigation methods for autonomous, high-speed vehicles in unstructured environments have demonstrated promising results in recent years. The success of fast-moving autonomous vehicles in off-road conditions can be attributed to two crucial factors. The first is learning-based vehicle control~\cite{williams_information_2017,lutter_deep_2019,  nagabandi_dexterous_2019,  yang_data_2019,shah2022gnm,kim_physics_2022}, and the second is the estimation of off-road terrain traversability~\cite{, kwak2017incremental_TE,ahtiainen2017normal,sevastopoulos2022survey,guan2022ga,contrastive_offroad,frey_2022locomotion,sathyamoorthy2022terrapn,TERP}.

Learning-based vehicle control has become crucial for autonomous vehicles since it facilitates optimal navigation performance while ensuring safety~\cite{kabzan_learning-based_2019, wu_daydreamer_2022, wang2021rough, kim_toast_2022}. Recent studies have demonstrated the effectiveness of combining learning-based vehicle dynamics models with model-based controllers~\cite{kahn2021badgr}. Also, estimating terrain traversability in off-road environments is essential for outdoor navigation. As off-road environments are rife with terrains that may generate uncertain vehicle reactions, failure to identify such uncertainty, which we refer to as \textit{perception failure}, could lead to navigation failure and even severe damage to the vehicle~\cite{Wellhausen_2020,pmlr-v155-ji21a,kerner2019novelty,ji2022proactive,seo2023learning}. Recent works present self-supervised methods for estimating traversability that take into vehicle-terrain interactions pertinent to navigation based on actual driving experiences~\cite{wellhausen2019should,Z_acoustic,Bekhti_verticala,Brooks_vibration,cai2022risk,gasparino2022wayfast,castro2022does}. They predict traversal costs derived from proprioceptive terrain interaction feedback using exteroceptive sensor measurements.


\begin{figure*}[t]
\centering
\includegraphics[width=0.9\textwidth]{Figures_Final/concept_compressed.pdf}
\caption{Diagram of our navigation method that minimizes control and perception uncertainty. While training the dynamics model, active exploration is conducted, which minimizes epistemic uncertainty using a Probabilistic Ensemble Neural Network~(PENN). Our traversability cost prediction network reduces aleatoric uncertainty through online adaptation using self-labeled data. Our method enables safe and efficient navigation in unstructured environments by combining uncertainty-minimizing control and perception methods.}
\label{fig:concept}
\vspace*{-0.2in}
\end{figure*}

However, the primary limitation of these learning-based methods of control and perception is that the uncertainty of the model significantly affects the dependability of predictions~\cite{nguyen-tuong_model_2011,palazzo2020domain,fan2021learning,cai2022risk,seo_scate_2023}. Insufficient data to learn vehicle dynamics may result in uncertainty and suboptimal control performance~\cite{kahn_uncertainty-aware_2017,parsi_active_2020}. Self-supervised traversability estimation methods have inherent uncertainty problems since the interaction data only provides supervision for actually traversed regions~\cite{Wellhausen_2020,seo_scate_2023,frey2023fast}. Moreover, since the terrain traversability is predicted from a limited sensor configuration (e.g., a sparse LiDAR point cloud)~\cite{fei2021pillarsegnet,cheng2021s3cnet,han2021planning,peng2022mass,shaban2022semantic}, the estimation entails substantial aleatoric uncertainty.

For the safe and reliable navigation of autonomous vehicles in unstructured environments, it is necessary to minimize the uncertainty of predictions regarding both control and perception. The uncertainty-aware navigation can be performed through data exploration during training or online adaptation during deployment~\cite{nagabandilearning, BanerjeeHarrisonEtAl2020,RichardsAzizanEtAl2021, visca2022Meta}. If uncertainty cannot be resolved, reliable navigation can be enforced by preventing vehicles from accessing uncertain state-action spaces or terrain with high levels of uncertainty.

This paper proposes a framework for autonomous vehicle navigation in unstructured environments. To this end, we employ an uncertainty-aware learning-based dynamics model and a self-supervised traversability estimation model. The vehicle dynamics model is trained through a model-based reinforcement learning framework that quantifies epistemic uncertainty. After learning the dynamics model, we generate an uncertainty-aware traversability cost map. Using the self-supervised traversability cost of terrains derived from interaction experiences, the traversability model is capable of online adaptation to minimize uncertainty. We integrate the dynamics model and traversability cost map with a sampling-based model predictive controller. The controller allows for safe navigation by optimizing the trajectory to avoid regions and actions with high uncertainty.
